\exnumber{09}
\extitle{UI I - Ex 10 - Défi UI - Spotify}
\exfiles{ex09/fichier Figma, nomme Name\_FirstName\_DeliverableName\_Date\_VersionNumber}
\exauthorize{Those named in the exercise}
\makeheaderfiles

\textbf{Instructions :}

Copiez cette image dans Figma
Spotify

\textbf{Livrable :}

\begin{itemize}
\item Un fichier Figma avec une copie de l'image et des éléments Figma
\item Please make sure you use:
\begin{itemize}
\item Linear color background
\item Groups
\item Named and organized layers
\item Apple design system (head bar)
\item Symbols and variants for the songs (1 page for the Symbols)
\end{itemize}
\item Déposez le fichier Figma sur Github, nommé ainsi : « Name\_FirstName\_DeliverableName\_Date\_VersionNumber ». Exemple : « Doe\_John\_MidFiPrototype\_05242024\_V1 ».
\item Ajoutez un titre dans le fichier.
\item Soignez l'aspect global du fichier. Il doit être clair, simple et lisible. Vous pouvez regarder canva.com, sans surcharger le dessin de détails.
\end{itemize}

\textbf{Pistes :}

\begin{itemize}
\item The font is GOTHAM
\item The frame is iPhone 13 Pro Max 428 x 926 px
\item Figma article about Symbols and Variants \url{https://help.figma.com/hc/en-us/articles/360056440594-Create-and-use-variants}
\end{itemize}

